\section{Comparison on root-preserving rotation orbits}\label{prof:local}

The next comparison transfers continuous-prior estimates to spins on
most orientations of each retained root orbit. Its essential quantitative
point is the signed cancellation between two spin and two Gaussian
replica integrals. An absolute bound on the matrix insertion would be too
large for the weighted covariance estimate.

The mixed cube--sphere tensor cancellation of
\citet[Lemma~3.5]{DuHuang2026CriticalOverlap} is the close predecessor of
this comparison. We retain the observation-dependent projection constraints
and a weighted signed matrix kernel throughout the calculation below.

\begin{proposition}[Covariance estimates on most root orientations]
\label{prof:median-covariance}
Let $\mathcal S$ be suitable base orbit data retained by the preceding
exclusions in \Cref{prof:orbits}. On $\mathcal S$, use the disintegration
$d\pi\,d\mathrm{Haar}$ before the determinant weight $w$.
Here $d\pi$ includes the scalar root factors and the imposed
rotation-invariant spectral restriction, $\mathrm{Haar}$ is normalized
measure on rotations fixing $\mathcal P$, and $w_0$ is the median
reference weight defined there.

Fix $1/2<\gamma<1$ and positive accuracies $\epsilon,\eta$, together
with the requested implied small constants below. Choose a sufficiently
small fixed rate $a_{\mathrm{orb}}>0$, then the upper bound on $r$
sufficiently small. At each scale $k=nr^3\ge\log^{10}n$, for sufficiently
large $n$, one may discard a further set $\mathcal E\subseteq\mathcal S$
satisfying
\[
 \int_{\mathcal E}w_0\,d\pi\le Cn^C e^{-a_{\mathrm{orb}}k}.
\]
On each remaining orbit, a set of Haar orientations of mass $>3/4$
satisfies simultaneously
\begin{equation}\label{prof:eq13}
\begin{split}
 \|\Sigma_s\|&\le (1+\epsilon)/s,\qquad
 \|A_0 \Sigma_s A_0\|\ll r^{-\gamma},\\
 \|m_s^{\mathrm{post}}-m\|&\ll r^{-1-\eta},\qquad
 \|A_0(m_s^{\mathrm{post}}-m)\|\ll r^{-\gamma/2}.
 \end{split}
\end{equation}
Here $m_s^{\mathrm{post}}$ and $\Sigma_s$ are the posterior mean and
covariance. Every implied small constant can be prescribed in advance by
shrinking the upper bound on $r$ and then taking $n$ sufficiently large.
The previously excluded unsuitable data retain their preceding
integrated bounds; $w_0$ is used only on $\mathcal S$. No bound uniform
conditional on each spectrum is asserted.
\end{proposition}

\begin{proof}
\par\smallskip\noindent\textit{Projection cuts and the three priors.}\enspace The plant is gauged to $\mathbf1$ as above. Let $\mathsf P$ be the
$n\times3$ matrix whose columns are the three orthogonalized profiles spanning
$\mathcal P$, normalized by $n^{-1}\mathsf P^T\mathsf P=I_3$.
Its first column is $m/\sqrt q$, and its second is the normalized vector
$\mathbf1-(M/q)m$; its third is the normalized residual of $y$ after
projection onto those first two columns. We write $\mathsf P_j$ for column
$j$, and empirical contractions of matrices always include the factor $n^{-1}$.

Put \(z=X-m,\ N=nq^3,\ L=\sqrt{n/q}\), and use the statistics
\[
 a=\sqrt n\,\delta,\quad \delta=\langle \mathsf P^T z\rangle=Q \vartheta,\quad Q=\operatorname{diag}(\sqrt q,q,q^{3/2}),\quad W=\sqrt N\, \vartheta.
\]
Here \(\vartheta\) denotes the scaled projection coordinate. For a pair write \(t=\langle(z-\mathsf P\delta)(z'-\mathsf P\delta')\rangle,\ U=n t/L\), so \(t=qU/\sqrt N\). The three-column notation always uses empirical contractions over sites. The exact length constraint is \(\langle z^2\rangle=b-2\sqrt q\,\delta_1\). We tilt by the common energy \(z^T(J-b)z/2\); with \(p\) the product of independent spins of means \(m\) (translated by \(-m\)), this gives precisely the posterior by the root equation.

We cut \(p\) smoothly to \(|\vartheta|\le\epsilon_0\) (one below
\(\epsilon_0/2\)), \(\epsilon_0>0\) an arbitrarily small fixed constant.
Except at a small rate in the true counting law with good spectrum, the
cut omits negligible normalized posterior mass (\(n^{-D}\), \(D\) any
fixed constant). Indeed a posterior spin can replace the plant in tests
with the equivariant axes \(m,y-lm\) by uniqueness and
\eqref{prof:eq9}, so
\(|\langle m z\rangle|\ll q,\ |\langle(y-lm)z\rangle|\ll q^3\).
The overlaps of plant and replica concentrate about \(q\) to any fixed
tolerance times \(q\) by \eqref{prof:eq1}, the spherical observation
estimates, \eqref{prof:eq4}–\eqref{prof:eq6} and \(q/r\simeq1\).

Thus \(|\langle z\rangle|\) (overlap with the gauged plant minus \(M\))
is bounded by arbitrarily small fixed times \(q\), and the subtraction
\((M/q)\langle mz\rangle\) in the \(e\) axis is likewise small.
Orthogonalizing the cubic axis uses \(|\langle y-lm\rangle|=O(q^3)\)
and residual norm asymptotic, in empirical units up to constant factors,
to \(q^{3/2}\), by the clipped Gaussian moment diagnostics and the
tail-square bound. Here multi-root fields do not affect counting claims:
the count of positive-sign regular roots is at most one plus the number
of negative-sign roots by degree, whose expectation on good spectrum is
exponentially small. Null singular fields may be ignored. The orbit
weight comparisons thus allow us to discard further orbits at a small
integrated rate measured by \(w_0\,d\pi\), so that the cut loss is
negligible on most Haar orientations.

Use two continuous auxiliary priors, not necessarily of mass one.
We fix their length-density units here. Put
\[
 E=\mathsf P/\sqrt n,\qquad
 \mathscr L(z)=\|z\|^2+2m^Tz,\qquad
 S_{\mathsf P}=E^TVE,\qquad \kappa_2=\langle v^2\rangle.
\]
For a Gaussian prior, the notation
\(\delta_{\rm D}(\mathscr L(z)-nb)\,dz\) means the coarea density
at the raw squared length \(nb\). Since this shell is
\(\|z+m\|^2=n\), it is surface measure divided by \(2\sqrt n\).
In particular, if \(T=(\mathscr L-nb)/L\), then
\(\delta_{\rm D}(\mathscr L-nb)=L^{-1}\delta_{\rm D}(T)\).
The spin prior already lies on this shell.

With \(\varphi_C\) the density of \(N(0,C)\), define the uncut
pre-energy measures
\[
\begin{split}
 g_0(dz)&=\mathfrak c_{g_0}\mathfrak r_{g_0}(a)
       \varphi_V(z)\delta_{\rm D}(\mathscr L(z)-nb)\,dz,\\
 o(dz)&=\mathfrak c_o\mathfrak r_o(a)
       \varphi_{bI}(z)\delta_{\rm D}(\mathscr L(z)-nb)\,dz,
\end{split}
\]
where
\[
 \mathfrak r_{g_0}(a)=e^{qa_1^2/\kappa_2},\qquad
 \mathfrak r_o(a)=
 e^{qa_1^2/b^2-\frac12a^T(S_{\mathsf P}^{-1}-b^{-1}I)a}.
\]
The constants \(\mathfrak c_{g_0},\mathfrak c_o\) are fixed by
the zero-projection saddle normalization below. Multiplication by
the common energy \(e^{z^T(J-bI)z/2}\) changes the underlying
Gaussian precisions to \(K\) and \(K_o\), respectively.

For \(Z\sim N(0,C_o)\), \(C_o=K_o^{-1}\), put
\(T_o(Z)=(\mathscr L(Z)-nb)/L\). The reference mass used
throughout is projection-uncut but still shell-constrained:
\[
 Z_*:=\int e^{z^T(J-bI)z/2}\,o(dz)
 =\mathfrak c_o\det(bK_o)^{-1/2}\frac{f_{T_o}(0)}{L}.
\]
Thus its normalized physical law is exactly the density
disintegration of \(N(0,C_o)\) conditional on \(T_o=0\).
Both factors in the last formula are invariant under rotations
fixing \(\mathcal P\), so \(Z_*\) is an orbit invariant.
We cut \(g_0\) by the same projection cutoff as \(p\), and multiply
it there by a projection-only factor to form \(g\). All cuts below
are separable across replicas.

The three priors serve different parts of the comparison. The reference
prior has an orbit-invariant tilted mass; the variable-variance Gaussian
has the precision needed for covariance control. The projection
multiplier aligns its single-replica marginal with the spin marginal.
What remains is a pair-density calculation in which the signed sum
cancels common errors and errors additive in the replica types.

\par\smallskip\noindent\textit{Single-prior projection saddles.}\enspace
We begin by defining the single-prior saddle functions. Write
\(p_\ell=\mathsf P_{\ell\bullet}^T\in\mathbb R^3\) for the
profile row at site \(\ell\). The independent site laws are
\[
\begin{split}
 \gamma_{p,\ell}&=\frac{1+m_\ell}{2}\delta_{1-m_\ell}
             +\frac{1-m_\ell}{2}\delta_{-1-m_\ell},\\
 \gamma_{g_0,\ell}&=N(0,v_\ell),\qquad
 \gamma_{o,\ell}=N(0,b).
\end{split}
\]
Their constraint vectors and per-site targets are
\[
\begin{array}{c|c|c|c}
 i&D_{i,\ell}(z)&c_i(\delta)&d_i=\dim D_{i,\ell}\\ \hline
 p&p_\ell z&\delta&3\\
 g_0,o&(p_\ell z,z^2)&(\delta,b-2\sqrt q\,\delta_1)&4.
\end{array}
\]
For a spin, \(z_\ell^2+2m_\ell z_\ell=v_\ell\) identically,
so its shell constraint is automatic; adding a square coordinate
would make the constraint covariance singular. For Gaussians the
square coordinate and its target are equivalent to the translated
length coordinate by a determinant-one triangular change.

Set
\[
 \Lambda_i(\theta)=\frac1n\sum_{\ell=1}^n
       \log\int e^{\theta^TD_{i,\ell}(z)}\,\gamma_{i,\ell}(dz).
\]
On the local interior branch let
\(\nabla\Lambda_i(\theta_i)=c_i(\delta)\) and
\(\mathsf H_i=\nabla^2\Lambda_i(\theta_i)\).
The branch and its positive Hessian are established on the small
projection cut below. With \(\mathfrak c_p=\mathfrak r_p=1\),
define
\[
 f_i^a(a)=\mathfrak c_i\mathfrak r_i(a)\,
 n^{3/2}(2\pi n)^{-d_i/2}(\det\mathsf H_i)^{-1/2}
 e^{n[\Lambda_i(\theta_i)-\theta_i^Tc_i(\delta)]}.
\]
This is a smooth saddle function per ordinary \(d^3a\), not an
exact density for an atomic spin projection. The factor \(n^{3/2}\)
converts the three raw projection sums \(\sqrt n\,a\) to \(a\);
the Gaussian raw length-density coordinate keeps its original
units. Primitive lattice covolumes are excluded here and will
belong to the mixed counting measures.

At zero projection the constraint Hessians are
\(S_{\mathsf P}\), \(\operatorname{diag}(S_{\mathsf P},2\kappa_2)\),
and \(\operatorname{diag}(bI_3,2b^2)\). Hence the choices
\[
 \mathfrak c_{g_0}=\sqrt{4\pi n\kappa_2},\qquad
 \mathfrak c_o=\sqrt{4\pi n b^5/\det S_{\mathsf P}}
\]
make all three saddle functions agree at zero. These are saddle
normalizations, not divisions by the exact shell masses.
The additional multiplier defining \(g\) on its cut is
\(f_p^a/f_{g_0}^a\).

These constant normalizations cost only polynomial for the unconditioned product Gaussian density integrals. On the cut we have
\begin{equation}\label{prof:eq14}
\log(f_i^a/f_o^a)=O(|\vartheta|+N |\vartheta|^3),
\end{equation}
also with the corresponding Taylor bounds after up to two scaled derivatives (bounded derivatives for the prefactor log, and \(O(N |\vartheta|^{3-j})\) for the rate-part log).

We include details of the expansions and bounds; all constants may use the scalar diagnostics above. In particular \(S_{\mathsf P}=I-O(q)\); \(Q^{-1}S_{\mathsf P} Q\) and its inverse are bounded. At typical sites, meaning the clipped region of those diagnostics with cutoff exponent sufficiently tiny, the profiles \(\mathsf P\) have bounded empirical absolute moments through any prescribed fixed degree, and \(y\) such moments at scale \(\sqrt q\). Outside that region \(\langle y^2{\bf1}_{\rm tail}\rangle=O(q^4)\); the scaled profiles \(\mathsf P Q\) are bounded there by \(C |y|\).

These bounds follow directly from \eqref{prof:eq9}, \(l=1+O(q)\), and the residual cubic normalization just noted. Centered moments involving \(\mathsf P\) and exponentially damped site fluctuations of order through fixed \(j\ge2\) thus have tail bounds for their profile factors up to \(C_j(1+q^{4-3j/2})\).

For \(p\) solve the linear-projection saddle by tilting by \(\mathsf P\theta\), \(\theta=O_Q(|\vartheta|)\), where \(O_Q\) denotes a bound in \(Q^{-1}\) units. Indeed on that parameter ball the response beyond linear \(v\mathsf P\theta\) has empirical \(L^2\) size \(O(q^{3/2}|Q^{-1}\theta|^2)\). Its scaled derivatives of order \(j\) through fixed orders have the bound \(C_j q^{3/2}|Q^{-1}\theta|^{\max(2-j,0)}\). This follows by the tanh derivatives (second bounded by a constant times the absolute small argument on typical sites, using the moment bounds), and on the tail the derivatives are exponentially damped since the new field differs from \(y\) by at most a small fraction of \(|y|\).

Thus scaled smooth inversion applies at uniform small radii. Evaluating the Legendre exponent at \(\theta=S_{\mathsf P}^{-1}\delta\) instead of the minimizer costs \(O(q^3 |\vartheta|^4)\) per site, by the ordinary Hessian bounds. The density exponent per site, apart from the Gaussian projection term \(-\delta^T S_{\mathsf P}^{-1}\delta/2\), is
\begin{equation}\label{prof:eq15}
-\sqrt q\,\delta_1^3-\delta_1^4/4+O(q^3|\vartheta|^3).
\end{equation}
Indeed the third log moment term is \(-\langle mv (\mathsf P\theta)^3\rangle/3\), with \(111\) profile average \(\langle mv \mathsf P_1^3\rangle=3\sqrt q+O(q^{3/2})\), \(112\) average \(O(q)\) by \eqref{prof:eq9}.

The fourth term is \(-\langle v(1-3m^2)(\mathsf P\theta)^4\rangle/12\), with \(1111\) average \(3+O(q)\), \(1112\) average \(O(\sqrt q)\). The fifth coefficient has a factor \(O(m)\), and the sixth remainder is bounded with tail damping. All other scaled monomials thus cost \(O(q^3)\) by typical moments and the tail bounds. In the two leading powers \(\theta_1\) can be replaced by \(\delta_1\). The estimates retain their Taylor derivative orders: the nonlinear Legendre correction does also by scaled smooth inversion and the stated derivative bounds (the unscaled shift in the minimizing coefficient has norm \(O(q^{3/2}|\vartheta|^2)\) with the corresponding derivatives).

For \(g_0\), first optimize the projection coefficients in the
formal saddle and then optimize the length coefficient. The
Gaussian projection exponent is exact. At this canonical projection
saddle, the length target differs from the mean by
\(-2\sqrt q\,\delta_1-\langle(v\mathsf P S_{\mathsf P}^{-1}\delta)^2\rangle\)
per site. This is a saddle mean, rather than the trace of an
exactly projection-conditioned finite-dimensional Gaussian.
The length curvature in the partial log moment is
\(2\langle v^2\rangle+O(q|\vartheta|^2)\). Thus the additional
exponent is the negative square of that deficit divided by
\(4\langle v^2\rangle\), to \(O(q^3|\vartheta|^3)\); its leading
quadratic cancels the explicit rank factor.

Since \(\langle(v\mathsf P S_{\mathsf P}^{-1}\delta)^2\rangle=|\delta|^2+O(q^3|\vartheta|^2)\), this gives exactly \eqref{prof:eq15}. These expansions use only the analytic bounded-variance Gaussian length tilt, adjusting the projection tilt by the corresponding three-dimensional linear solve. Length coefficients are \(O(q |\vartheta|)\), and scaled Taylor derivative estimates are unchanged. The calculation for \(o\) replaces the variance profile by constant \(b\), giving the same terms. The Gaussian central prefactors for all three priors have uniformly bounded log derivatives in \(\vartheta\) by the nonsingular saddle Hessians and the profile bounds. This proves \eqref{prof:eq14}.

\par\smallskip\noindent\textit{The signed pair comparison.}\enspace
We next form the pair saddle functions on
\(|\vartheta|,|\vartheta'|\le\epsilon_0\) and
\(|t|/q\le\epsilon_1\), where \(\epsilon_1\) is fixed small
and \(\epsilon_0\) may be much smaller. For types
\(i,j\in\{p,g_0,o\}\), put
\[
\begin{split}
 D_\ell^{ij}(z,z')&=(D_{i,\ell}(z),D_{j,\ell}(z')),\\
 c^{ij}&=(c_i(\delta),c_j(\delta')),\\
 h_\ell(z,z')&=(z-p_\ell^T\delta)(z'-p_\ell^T\delta'),
\end{split}
\]
so \(n^{-1}\sum_\ell h_\ell=t\) at the projection constraints.
Define
\[
 \Lambda_{ij}(\theta,\beta)=\frac1n\sum_\ell
 \log\iint e^{\theta^TD_\ell^{ij}+\beta h_\ell}
       \,\gamma_{i,\ell}(dz)\gamma_{j,\ell}(dz').
\]
Keep the targets and the \(\delta\)'s inside \(h_\ell\) fixed
when differentiating. On the local interior branch let
\(\nabla_\theta\Lambda_{ij}(\theta_\beta,\beta)=c^{ij}\), and put
\[
 F_{ij}(\beta)=
 \Lambda_{ij}(\theta_\beta,\beta)-\theta_\beta^Tc^{ij}
 -\bigl[\Lambda_{ij}(\theta_0,0)-\theta_0^Tc^{ij}\bigr].
\]
Thus \(F_{ij}(0)=0\); it is the partial infimum over the
single-constraint coefficients, minus its value at zero cross
tilt. It is a formal optimized log moment, not the conditional
log moment of an unsmoothed atomic event.

At the optimizing law write
\[
 D_2=\langle\operatorname{Cov}(D_\ell^{ij},D_\ell^{ij})\rangle,
 \quad v_\beta=\langle\operatorname{Cov}(D_\ell^{ij},h_\ell)\rangle,
 \quad w_\beta=\langle\operatorname{Var}(h_\ell)\rangle.
\]
These are averages of sitewise covariances. Implicit differentiation
and the envelope identity give
\[
 \theta_\beta'=-D_2^{-1}v_\beta,\qquad
 F_{ij}'(\beta)=\langle\E_\beta h_\ell\rangle,\qquad
 F_{ij}''(\beta)=w_\beta-v_\beta^TD_2^{-1}v_\beta.
\]
In particular the last expression is the variance after regressing
the cross statistic on the single constraints.

Let \(\beta_{ij}(t)\) solve \(F_{ij}'(\beta_{ij})=t\), with
\(t=LU/n\). The formal joint function in ordinary
\(d^3a\,d^3a'\,dU\), retaining every raw Gaussian length density, is
\[
\begin{split}
 \widehat f_{ij}^{\,a,a',U}
 ={}&\mathfrak c_i\mathfrak c_j
       \mathfrak r_i(a)\mathfrak r_j(a')\,
       n^3L(2\pi n)^{-(d_i+d_j+1)/2}\\
 &\times[\det D_2^{ij}(\beta_{ij})F_{ij}''(\beta_{ij})]^{-1/2}
 e^{n[\Lambda_{ij}(\theta_{\beta_{ij}},\beta_{ij})
        -\theta_{\beta_{ij}}^Tc^{ij}-\beta_{ij}t]}.
\end{split}
\]
The factor \(n^3L\) converts the six raw projection sums and the
raw cross sum \(nt=LU\). Dividing this formula by its two single
saddle functions, and converting \(U\) back to \(t\), gives
\[
 \widehat h_{ij}^{\,t}(t\mid a,a')
 =\sqrt{\frac n{2\pi}}
 \left[\frac{\det D_2^{ij}(0)}
 {\det D_2^{ij}(\beta_{ij})F_{ij}''(\beta_{ij})}\right]^{1/2}
 e^{n[F_{ij}(\beta_{ij})-\beta_{ij}t]}.
\]
The determinant quotient uses each pair's own constraint dimension:
\(6,7,8\) for \(pp\), mixed pairs, and two Gaussian replicas.
The Gaussian shell units and all explicit rank factors cancel in
this quotient. Define the dimensionless formal ratio
\[
 \mathcal R_{ij}=
       \widehat h_{ij}^{\,t}/\widehat h_{oo}^{\,t}.
\]
It is equally the ratio in \(U\) units, since both densities then
acquire the factor \(L/n\). The full formal pair ratio is
\[
 \frac{\widehat f_{ij}^{\,a,a',U}}
      {\widehat f_{oo}^{\,a,a',U}}
 =\frac{f_i^a(a)f_j^a(a')}{f_o^a(a)f_o^a(a')}\mathcal R_{ij}.
\]
This is an algebraic quotient of saddle formulas. The local-limit
argument below relates it to the actual mixed statistic measures.
After the multiplier defining \(g\), both effective single factors
are \(f_p^a\); the projection cut does not change the conditional
calculation.

The signed comparison and its verification below use the four ordered
pairs \(i,j\in\{p,g_0\}\). With signs \(s_p=1,s_{g_0}=-1\) and
\(E_*=1+|W|+|W'|+|U|\), the full signed pair-density comparison needed is
\begin{equation}\label{prof:eq16}
\frac{f_p^a(a)f_p^a(a')}{f_o^a(a)f_o^a(a')}
 \sum_{i,j\in\{p,g_0\}}s_i s_j\,\mathcal R_{ij}
 =
 \frac{q}{\sqrt N}\,e^{c_2 q U^2}(c_1 U+c_3' U^3)
 +O\!\left(\frac qN E_*^C e^{\zeta E_*^2}\right).
\end{equation}
Here cuts can multiply both sides; the coefficients are bounded and independent of the statistics. The constant \(\zeta\) can be made as small as needed by reducing the projection cut and the upper \(r\). Central-prefactor relative errors \(O(q/N)\) are allowed.

Here is a verification. Abbreviate \(F_{ij}\) by \(F\) for the
pair under discussion. At \(\beta=0\) the canonical laws are
independent single-replica laws. Write their site means and
variances as \(\lambda_i,\sigma_i\), where \(\sigma_i\) denotes
variance rather than its square root, and put
\(d_i^0=\lambda_i-\mathsf P\delta\), with the primed target for
the second replica.

Uniformly with scaled \(\vartheta\) derivatives,
\begin{equation}\label{prof:eq17}
\begin{split}
 \|d_i^0-(v\mathsf P S_{\mathsf P}^{-1}\delta-\mathsf P\delta)\|_{2,\mathrm{emp}}
     &\le C q^{3/2}|\vartheta|^2,\qquad
 \|d_i^0\|_{2,\mathrm{emp}}\le C q^{3/2}|\vartheta|,\\
 \|\sigma_i-\langle\sigma_i\rangle\|_{2,\mathrm{emp}}&\le Cq,\qquad
 \langle\sigma_i\rangle=B(\vartheta)-\|d_i^0\|_{2,\mathrm{emp}}^2,\quad
 B(\vartheta)=b-2\sqrt q\,\delta_1-|\delta|^2 .
 \end{split}
\end{equation}
Indeed the common linear response and nonlinear bounds were just obtained (for Gaussian the shifted length coefficient changes the response only nonlinearly).

The variance assertions use tanh or Gaussian differentiation as above; the last identity uses the exact target, automatic also for spins. The first bounds hold in any needed fixed typical-site empirical moment; on the tail \(d_i^0\) and its scaled derivatives are bounded by \(C |y|\).

Consequently
\[
 F'(0)=\langle d_i^0 d_j^0\rangle,\qquad
 F''(0)=\langle\sigma_i\sigma_j\rangle+O(q^3(|\vartheta|+|\vartheta'|)^2).
\]
In the conditional curvature the centered-fluctuation product is orthogonal to the constraints, and the two linear remainder terms are regressed on them, proving the second formula, also to differentiated orders. For the size bounds with scaled projection derivatives used here, \(\sigma_i=v\) at zero, and \(\sigma_i\) and its scaled derivatives after empirical mean subtraction have \(L^2\) size \(O(q)\): the logistic first derivative of variance brings a factor of the magnetization, and higher Taylor terms in the site shift or the Gaussian length coefficient already pay order \(q\) (use the same parameter estimates as for the single saddles).

On the tail all changes of the site fields and their scaled derivatives cost \(O(1+|y|)\) besides bounded quadratic coefficients, against fluctuation moments damped exponentially. Constraint covariance matrices and their needed scaled derivatives at the independent saddles are bounded, with eigenvalues bounded below using just bounded-profile sites. Regression coefficients for the two linear remainders and their derivatives therefore use ordinary unscaled constraint units (profiles \(\mathsf P\), squares for Gaussian), with size \(O(q^{3/2})\) by \eqref{prof:eq17}.

Thus \(F'(0)\) to quadratic order, of size \(O(q^3(|\vartheta|+|\vartheta'|)^2)\), is common to all four pairs. The curvature discrepancy from \(B(\vartheta)B(\vartheta')\) and its derivatives are \(O(q^2)\); it is common at zero projections, and to linear order varies by an \emph{additive} sum of single-type terms. The third derivative is \(O(q)\) at \(\beta=0\), with \(O(q)\) projection derivatives. Indeed it is the third cumulant averaged over sites of the regression-adjusted cross observable.

At zero cross tilt the product of centered site fluctuations contributes the average product of individual third moments, \(O(q)\) with derivatives. The linear corrections have coefficients \(d_i^0,d_j^0\), and regression coefficients on site constraints \(D=(\mathsf P z,\mathsf P z',\text{applicable squares})\) of size \(O(q^{3/2}(|\vartheta|+|\vartheta'|))\). Their cumulant contributions and derivatives satisfy the stated bound by typical moments in \eqref{prof:eq17}; on tails they use at most exponentially damped polynomials in \(|y|+1\), of tail average \(O(q^3)\).

For clarity, fourth cross derivatives stay bounded up to \(|\beta|\le C\epsilon_1 q\). The constraint-coefficient shifts have first derivatives of norm \(O(q^{3/2}(|\vartheta|+|\vartheta'|)+|\beta|)\), second derivatives bounded. Indeed bootstrap their first derivatives at \(O(q)\), so shifts are \(O(q^2)\), giving only small parameter changes even in \(Q^{-1}\) units. The constraint Hessian \(D_2\) (averaged covariance of \(D\)) stays well-conditioned. Centered site fluctuations of \(z,z'\) remain exponentially damped in \(|y|\) on tails, with bounded means: for a coupled Gaussian integrate the small quadratic tilt exactly, and for spins the conditional odds or Gaussian-integrated odds still have fields \(y+O(1)+o(1)|y|\).

Thus the regression-adjusted observable \(H=(z-\mathsf P\delta)(z'-\mathsf P\delta')+\theta'_\beta D\), with \(\theta_\beta\) the optimizing single coefficients, has bounded empirical centered fourth moment, using the \(O(q)\) coefficient times the worst \(\mathsf P_3\) profile. Implicit differentiation then bounds \(\theta''_\beta\) by \(C\|\langle\kappa(D,H,H)\rangle\|\le C'\), closing the bootstrap (or use continuation at a larger threshold and small \(\epsilon_1\)); \(\kappa\) denotes joint cumulant.

The fourth derivative of \(F\) uses only the fourth cumulant of \(H\) and the constraint regression of its square. Similarly \(\partial_\beta D_2\) is bounded (using at most two constraint factors and one \(H\)), and \(\partial_\beta^2 D_2=O(1+q^{-1/2})\): with two \(H\)'s the worst four-profile tail costs \(Cq^2 q^{4-6}\), while the \(\theta_\beta''\) term costs a third constraint moment. At zero cross tilt \(\partial_\beta D_2\) varies with bounded scaled projection derivatives by the same estimates. Hessian lower bounds throughout use just typical sites; in particular \(F''\) stays bounded below. All these statements apply to mixed constraint dimensions without dummy square variables for spins.

The pair conditional exponent is therefore
\[
 n\left[-\frac{(t-F'(0))^2}{2F''(0)}
   +\frac{F'''(0)(t-F'(0))^3}{6F''(0)^3}+O((t-F'(0))^4)\right].
\]
The preceding calculations give the cross mean, curvature, and third
cumulant at the independent saddle, with their projection derivatives.
We now convert them into a conditional density. Both the prefactor and
the Legendre exponent must be retained, because each contributes to the
odd term in the signed comparison.

For the unweighted reference, the exact conditional law uses two
independent uniform directions in dimension \(n-3\). Writing
\(B=B(\vartheta)\), \(B'=B(\vartheta')\), its density is
\[
 h_{oo}^{t,\mathrm{ex}}(t\mid a,a')
 =\frac{c_{n-3}}{\sqrt{BB'}}
       \left(1-\frac{t^2}{BB'}\right)^{(n-6)/2},
 \qquad
 c_d=\frac{\Gamma(d/2)}{\sqrt\pi\,\Gamma((d-1)/2)}
\]
on \(|t|<\sqrt{BB'}\). The formal \(\widehat h_{oo}^{\,t}\)
has the same power and replaces \(c_{n-3}\) by
\(\sqrt{n/(2\pi)}\); their ratio is \(1+O(1/n)\), a permitted
relative error. The reference Legendre rate exponent is
\(\frac n2\log(1-t^2/(BB'))=-nt^2/(2BB')+O(nt^4)\).
The remaining factor \((1-t^2/(BB'))^{-3}\) belongs to the
reference saddle prefactor. The pair conditional prefactor is
\((F''(\beta)\det D_2(\beta)/\det D_2(0))^{-1/2}\) at the dual
\(\beta\), while the reference prefactor is
\((BB')^{-1/2}\) up to \(O(t^2+1/n)\) relative error.
The curvature part at \(\beta=0\) thus has a common constant,
additive linear log terms of size
\(O(q^2(|\vartheta|+|\vartheta'|))\), and error
\(O(q^2(|\vartheta|+|\vartheta'|)^2)\). The log change at
\(\beta\) is a projection-independent linear term in \(t\),
up to \(O(qE_*^C/N)\), by the preceding derivative bounds and
dual inversion.

We organize the exponent difference by its dependence on the pair
types. Common terms and terms additive in one type cancel after the
signed sum. The remaining leading terms are odd in the normalized
overlap; quadratic products of the smaller terms fit the remainder.

In the exponent difference, the \(t^2\) terms give a common \(c_2 q U^2\) plus additive single-type terms \(O(q U^2 (|W|+|W'|)/\sqrt N)\) and error \(O(q E_*^C/N)\). The term linear in \(F'(0)\) is common up to that error, of size \(O(q E_*^3/\sqrt N)\); its square contribution to the rate is negligible. The cubic term at leading order is odd, of order \(qU^3/\sqrt N\), with coefficient at zero projections. All fourth-order remainders cost \(Cq E_*^C/N\). Indeed \(nt^2=q^{-1}U^2,\ nt^3=U^3/\sqrt N,\ nt^4=q U^4/N\).

In \(F''(0)-B B'\) the terms through projection degree one are just from the empirical covariance of the two site-variance profiles (both \(v\) at zero). Hence the inverse-curvature discrepancy has a common constant and single-replica-summed linear terms, with remainder \(O(q^2 E_*^2/N)\). In \(ntF'(0)\), the response product has a common bilinear part, leaving error \(O(q E_*^C/N)\) after multiplication. The log-prefactor shift uses \(\beta=(t-F'(0))/F''(0)+O((t-F'(0))^2)\); even its \(O((1+q^{-1/2})\beta^2)\) Taylor remainder fits.

Products of two of the small terms cost \(O(q E_*^C/N)\). Exponentiating (with the small Gaussian-growth envelope stated, using small cutoffs), the common and additive terms cancel from the signed sum, leaving just the claimed odd terms. Finally \eqref{prof:eq14} allows the outside projection factor to be replaced by one on those odd terms to the stated remainder. We check the remainder sizes term by term, including at the ends of
the cut window. Writing \(\mathcal W=|W|+|W'|\), the linear-mean term \(ntF'(0)/F''(0)\) has common leading polynomial of size \(C q |U|\mathcal W^2/\sqrt N\) and remainder \(C q |U|\mathcal W^3/N\), while \(nF'(0)^2\le C q^3\mathcal W^4/N\).

The inverse curvature difference from \(1/(B(\vartheta)B(\vartheta'))\) has quadratic-and-higher Taylor error \(C q^2\mathcal W^2/N\). The cubic coefficient discrepancy from its value at zero projections is \(Cq\mathcal W/\sqrt N\), and replacing \(t-F'(0)\) by \(t\) inside that cube also fits the stated error. Hence besides the common \(c_2 q U^2\) and projection-independent log-prefactor constant (common because both variances are \(v\) at zero), log contributions retained are of size \(C q E_*^3/\sqrt N\), common, single-type summed, or odd as indicated.

In exponentiating, both the leading polynomials in the log and the full log deviations (using \(n|t-F'(0)|^4\) for the Legendre remainder) have absolute upper bounds by constants plus \(Cq E_*^2\) on these cuts, before any outside projection log which adds \(O(|\vartheta|+|\vartheta'|+(|\vartheta|+|\vartheta'|)E_*^2)\). This yields the stated Gaussian-growth envelope even at window ends, and the common and single-type-summed terms indeed drop out after multiplying by the common projection weights and summing both signs.

The zero-projection coefficients and the exact common/additive
cancellation are also written explicitly in
Lemma~\ref{prof:surviving-coefficients}. This proves the smooth expansion
\eqref{prof:eq16}. Individually the full ratios have at most the same Gaussian-growth envelope, and differ from one on bounded normalized statistics by \(O(q)+o(1)\).

\par\smallskip\noindent\textit{Density inversion, tails, and the common angular kernel.}\enspace
The four pair laws have different statistic measures. We keep their
lattice coordinates and translated lengths exact, and smooth only the
other coordinates. Each remainder is estimated on its own measure.
The smooth leading terms are transferred to common volume before we
perform the signed cancellation.

\begin{itemize}
\item On the central region just used, one can evaluate orbit integrals as lattice sums and smooth densities with the saddle formulas, each to relative remainder \(O(q/N)\). Use lattice variables \(\sum x\) for each sign replica and \(\sum xx'\) for a pair of signs, on the translated lattice generated by differences of sign patterns. One can use \(a_1,a_3\) as the remaining continuous-approximation projection coordinates for sign replicas. Smooth all nonlattice variables at inverse arbitrarily high polynomial scale with compactly supported smooth kernels, \emph{except} exact Gaussian length constraints \(\sum(z^2+2mz)=nb\).

Changing an angular test by this smoothing costs an arbitrarily small inverse polynomial absolute error times the corresponding positive partition integrals (on slightly enlarged regions). Indeed projections and residual Gram factor parametrize the complement rotations, with Gram nonsingular here, so energies, pre-factors (in log), smooth cutoffs and the bounded-degree polynomial tests have polynomial variation. Only the Gaussian length itself needs remain imposed on the original site integral; inverse site variances need not be bounded.

Projection multipliers also have polynomial log derivatives. When lattice variables are present keep them exact in this procedure. For instance for two signs one may hold \((\sum x,\sum x',\sum xx')\) and smooth the two nonconstant projection coordinates for each sign; \(a_2,a'_2,U\) are recovered using those constant-coordinate sums, the other \(a\)'s and known site data. These are nonsingular affine or triangular coordinate changes (equivalently use the three raw statistics in site units together with the nonconstant profiles).

For a mixed pair the raw lattice variable is just \(\sum x\) for the sign, and one also smooths the nonlattice overlap. Deterministic statistic Jacobians are understood in saddle densities and in converting lattice mass units. The tilt optimization using projected deviations rather than raw cross products is an affine change at the fixed projection constraints, with the same partial Legendre rate.

\end{itemize}
For the density estimation, tilt to the canonical product saddle at the evaluation point; variation of tilting factors across the mollifier costs negligible relative error. Fourier inversion there gives the Gaussian prefactor, including lattice covolume, to relative \(O((1+q^{-2})/n)\). Indeed covariance in site-sum units (normalized by \(\sqrt n\)) is bounded positive definite. We can use the well-conditioned equivalent raw coordinates (including \(zz'\) or \(xx'\) rather than the projection-subtracted statistic).

Maximum centered site variance in any such direction is \(O(q^{-3})\), average third absolute and fourth moments \(O(1+q^{-1/2})\), \(O(1+q^{-2})\), by the bounds above. In the central Fourier expansion up to small times \(\sqrt N\) in normalized frequency units, each site factor stays close to one. The cubic log term integrates to zero at first order by oddness, and the fourth remainder and squared cubic term, with Gaussian damping, give the indicated error. One can evaluate at the exact mean here even with smoothing (or keep its transform, which changes central frequencies negligibly).

For more detail about Fourier tails, bounded typical sites damp through small constant frequencies in unnormalized (single-site) units by nonsingular covariance on that subset. Up to arbitrarily high polynomial frequencies thereafter there is exponential-in-size damping off the central region of the reciprocal-lattice fundamental cell. For a sign pair, phase comparisons on the four patterns test the two nonconstant profile combinations plus constants by single-spin flips (with both signs for the other spin).

Thus if nonconstant coefficients for either replica are not tiny, the nonlattice diagnostics give damping; with tiny coefficients on the bounded typical sites, only the reciprocal lattice of the constant three-pattern statistics can escape the test. More explicitly these statistics are the coefficients of \(x,x',xx'\), and equality of the four phases characterizes exactly that dual lattice. Typical-site pattern probabilities are bounded below. The single sign test works identically. For Gaussian variables on typical sites condition on any site signs first.

Quadratic Gaussian phases separated from zero give contraction irrespective of the linear part. With tiny quadratic coefficients, Gaussian linear phases separated from zero on a fixed fraction also give contraction; for a mixed sign pair one tests both sign choices to detect the cross coefficient. The Gaussian profile columns are well-conditioned on the bounded typical sites. Their conditional real means are bounded there. Thus failure of contraction from these tests requires all Gaussian-related coefficients tiny, when on this subset they only perturb the sign phase test by a small amount.

This argument can take successively small fixed cutoffs, or pass to sequences where contraction vanishes (quadratic phases must tend to zero, then Gaussian linears for each sign choice, then the empirical sign-phase dispersions). Beyond polynomial frequencies the mollifiers suffice; in unsmoothed length directions the Gaussian quadratic factors on typical sites themselves give a bound \((1+c|\omega_{\rm lengths}|^2)^{-c'n}\), also in the joint inversion. Here length direction means the single-site quadratic coefficients; the linear parts do not affect this bound.

Thus one is approximating ordinary joint constraint densities in the Gaussian length dimensions, and possibly smoothed/lattice joint masses elsewhere. Factoring a result into single prefactors and a conditional saddle expression \eqref{prof:eq16} just uses determinant Schur factorization and rate subtraction at the level of these formulas; no division by tiny unsmoothed atomic projection probabilities is required. These estimates justify inversions in the different constraint dimensions and also the denominator calculations for conditional densities.

\begin{itemize}
\item Overlap tails outside \(|t|\le\epsilon_1 q\), with the projection cuts, are negligible (\(e^{-cN}\) times polynomial) in each pair orbit integral relative to \(Z_*^2\), even for polynomial tests. First take \(\epsilon_1 q<|t|<c_*\) with \(c_*\) small fixed. Conditional density ratios can here be estimated as upper bounds on polynomially small boxes. At independent single canonical laws use exponential Markov on \((z-\mathsf P\delta)(z'-\mathsf P\delta')\) up to tilts \(C c_*\), without having to keep its constraint regressions optimized.

Its mean is \(O(q^3)\); unconditional curvature differs from \(B B'\) by \(O(q^2)\) by \eqref{prof:eq17}; third cumulant is \(O(q)\), fourth derivatives bounded. Indeed single saddle parameters remain fixed in this tilt and the damping on tails is preserved (no large new profile coefficients, \(|\mathsf P\delta|\le C\epsilon_0 |y|\) there, and site constants do not enter higher cumulants). Thus relative to the reference conditional overlap density the log cost, besides polynomial losses, is at most
\end{itemize}
\[
 C n(q^3|t|+q^2t^2+q|t|^3+t^4).
\]
Imposing exact Gaussian length densities costs at most polynomial under these tilts by bounded typical-site Gaussian smoothing as in the inversion bound.

Projection marginals relative to reference cost at most \(\exp(C\epsilon_0^3 N)\) besides constants by \eqref{prof:eq14}, including our modification. Boxes can straddle lattice coordinates without needing any matching lower bound there for the numerator. In the physically tilted reference \(C_o\), the length deficit of one of \((z\pm z')/\sqrt2\) is at least \(n|t|/2\) from trace, on taking \(\epsilon_0\ll\epsilon_1\). This costs \(\exp(-c n |t|^3)\) by negative square tilt of strength small-times \(t^2\) (and conditioning at polynomial cost).

Indeed the covariance loses normalized trace at most \(C\sqrt{\xi}+o(q)\) at added precision \(\xi\ge0\) small by \eqref{prof:eq1} and finite-rank comparison; lengths have central denominators and upper density bounds polynomial under the tests by Gaussian Fourier damping. Equivalently exact shells here can use buffered radial conditioning on the spheres translated by \(m\).
The absorption is uniform down to the moving endpoint
\(|t|=\epsilon_1q\), since
\[
 C\left(\frac{q^3}{t^2}+\frac{q^2}{|t|}+q+|t|\right)
 \le C\left(\frac q{\epsilon_1^2}+\frac q{\epsilon_1}+q+c_*\right)
\]
is the scalar cost divided by \(n|t|^3\), while the projection cost
has relative size at most \(C(\epsilon_0/\epsilon_1)^3\).
Hold \(\epsilon_1\) fixed, choose \(c_*\) and then the upper bound
on \(q\) small, and finally choose \(\epsilon_0/\epsilon_1\) small.
The reference penalty then leaves a fixed multiple of \(n|t|^3\),
which is at least a fixed multiple of \(N\). Angular comparisons
over these boxes have polynomially controlled variation.

For larger \(|t|\), use directly the angular semicircle pair upper bound from \eqref{prof:eq6} in the complement of \(\mathcal P\), whose Gram eigenvalues here are \(1\pm t+o_{\rm small}(1)\), and ignore projected components at \(o_{\rm small}(n)\) cost. The log gain besides independent baselines is at most \(nt^2/2+o_{\rm small}(n)\). Indeed \(\log Z_*/n\ge1/4-1/2-o_{\rm small}(1)\) from its \(C_o\) formula by central length density, the semicircle log determinant upper bound and \eqref{prof:eq10}; finite-rank and prior normalization changes do not increase the log determinant by a nonvanishing order per spin.

The sign-sign prior pays the strict Cramér slack already used, with biased signs and translations changing bounds by \(o_{\rm small}(n)\). For Gaussian pairs, simultaneous length/overlap Chernoff gives at least the rate \(-\frac12\log(1-t^2)\) in the small-\(q\) limit, strictly larger than \(t^2/2\) throughout this range. For two Gaussian types this uses fixed tilts of the Gram matrix, with log moments tending to the iid unit-variance values (\(v\le1\)); for mixed types condition on the sign vector, of bounded entries after translation and near-unit empirical length, and tilt the one Gaussian length and its projection on that vector, with the same limiting rate.

This follows immediately from the normal moment formulas since \(\langle|v-1|\rangle=q\); the mixed vector entries are bounded. Interior finite tilts suffice including by approaching the endpoints, using compactness and only requesting strict slack. Length targets differ from 1 negligibly, and their exact Gaussian densities again cost polynomial under the tilts. Projection/pre-rank factors cost \(o_{\rm small}(n)\). These estimates also handle singular pair Grams directly, by splitting the fixed overlap range into sufficiently small fixed bins. This proves the tail assertion.

\begin{itemize}
\item Here is the accuracy of kernel integration and of converting lattice sums to ordinary integrals. Under the tilted reference, condition two \(C_o\)-normals independently on their length shells, whose density denominators in \(L\) units are bounded above and below as in \eqref{prof:eq2}. Conditional on projections \(a,a'\) before length conditioning, their white coordinates \(\eta_i\) have fixed projections described by standard-normal coordinates \(\alpha,\alpha'\) (exactly whitened \(a,a'\)), and fresh standard orthogonal complements.

In particular \(|W|\le C|\alpha|\), and raw projection covariance eigenvalues are bounded below since \(K_o\) has bounded norm. Lengths and \(U\) are linear-quadratic forms in the remaining normals. The quadratic parts in a Fourier direction of norm one have operator and Hilbert-Schmidt norms at most \(C/\sqrt N,C\), with spread giving determinant decay \((1+c |v_f|^2/N)^{-c'N}\) at frequency \(v_f\). Indeed they use \(C_o/L\) tensored with a two-by-two symmetric matrix (whose norm in such a direction is bounded below), up to the finite-rank compressions/subtractions, and \eqref{prof:eq10} applies.

The length linear part uses \(\|C_o^{1/2}m\|/L\le C\). Consequently central densities with any fixed derivatives in \(a,a',U\), including the projection density, are bounded by
\end{itemize}
\begin{equation}\label{prof:eq18}
C(1+|\alpha|+|\alpha'|+|U|)^C e^{-c(|\alpha|^2+|\alpha'|^2+U^2)}
 \big|\det(\partial\alpha/\partial a)\big|^2 .
\end{equation}
To see the \(U\)-tail factor before inversion, tilt in \(U\) by parameter small-times the target \(U\); up to \(\epsilon_1\sqrt N\) the quadratic norm shift is small. Its centered determinant cost is \(\exp(C\,{\rm parameter}^2)\) (off-diagonal tilt), and the projected translations cost at most \(\exp(o_{\rm small}(1)(|\alpha|^2+|\alpha'|^2))\). The tilted spread is unchanged up to constants.

For the weighted matrix test, subtract $C_o$ times the mass before
forming the Hilbert--Schmidt square. The resulting insertion must be
averaged with its sign: its absolute value would lose the Gaussian
cancellation at the scale required here. Its insertion is
\[
 H=(\eta_1^TD\eta_2)^2-\eta_1^TD^2\eta_1-
   \eta_2^TD^2\eta_2+\operatorname{Tr}D^2,
 \qquad D=C_o^{1/2}A_0^2C_o^{1/2},\quad \|D\|\le C.
\]
The next kernel estimate controls this signed average and, after the
projection variables are integrated out, its smaller odd part.

Lemma~\ref{prof:gaussian-signed-kernel}, applied with $C=C_o$,
$A=A_0$ and $E=\mathsf P/\sqrt n$, proves the signed-density bound
\eqref{prof:eq18} times $Cnq$, including its fixed derivatives.
We verify its hypotheses explicitly. The formula for $K_o$,
$S_{\mathsf P}=I-O(q)$, and $\|J\|\le3$ give $\|K_o\|\le C$;
therefore $C_o\succeq C^{-1}I$. The lower precision bound,
trace-square bound, and edge spread follow from \eqref{prof:eq10}.
Its first projection estimate also gives
\[
 \|C_o^{1/2}m\|^2=nq\,u^TC_ou\le Cn/q=CL^2.
\]
Thus
\[
 \operatorname{Var}(\|z\|^2+2m^Tz)
 =2\operatorname{Tr}C_o^2+4m^TC_om\asymp L^2,
 \qquad \operatorname{Tr}C_o=nb+o(L).
\]
The lower variance bound uses the edge spread. The density estimate
\eqref{prof:eq2}, with its Fourier upper bound, now puts both shell
density denominators between fixed positive constants in $L$ units.
After integrating out $a,a'$, the odd parts of both the mass and matrix
densities gain $N^{-1/2}$.

The lemma proves this by comparing the
cross-frequency signs before projection conditioning; hence it does
not assume an artificial symmetry of the translated shells.
Projection cuts multiplying the leading odd term can be removed at
exponentially small cost, since \eqref{prof:eq18} gives Gaussian tails
in the whitened projection variables.

Lattice sums after the density approximation use steps bounded by \(C/\sqrt n\) in the raw \(a\)'s and \(C/L\) in full-overlap \(L\) units (fixing the other projection coordinates first, with affine shifts). Their covolumes are exactly the prefactor ones from Fourier inversion. Passing to \(U\) has bounded derivatives on the cut: besides subtracting \(a\cdot a'/L\), converting a full \(x\)-overlap to a \(z\)-overlap subtracts the \(m\)-linears with derivative \(\sqrt{nq}/L=q\). Converting constant-axis sums to \(a_2\)'s uses \(\mathsf P_2\)'s formula with bounded coefficients.

Thus the centered-cell integration error for smooth leading terms is \(O(1/n)\) times their respective polynomial-moment envelope scale. More explicitly use the individual expansions leading to \eqref{prof:eq16} with second derivatives of the explicit constant/common/additive and odd terms, keeping the projection factors via \eqref{prof:eq14}; in raw \(a,a',U\) units these derivatives cost just polynomials times \(e^{\zeta E_*^2}\). Remainders \(O(q E_*^C e^{\zeta E_*^2}/N)\) require no derivatives.

The type of error matters here. Density-ratio errors, including relative
saddle errors, are scalar functions of the statistics and multiply the
\emph{same angular kernel} as the leading terms. We make the exact
denominator in this assertion explicit. Use the chart \(\Phi\),
Jacobian \(J_\Phi\), mixed measures \(dm_i\), and exact unweighted
reference density \(f_{oo}^{\rm unw}\) from
Lemma~\ref{prof:common-angular-kernel}. For the reference single
projection, isotropy gives the exact formula
\[
 f_o^{a,\rm ex}(a)=\mathfrak c_o\mathfrak r_o(a)
     \varphi_{bI_3}(a)f_{b\chi^2_{n-3}}(\rho(a)^2),
 \qquad \rho(a)^2=nb-2\sqrt{nq}\,a_1-|a|^2.
\]
Its ratio to \(f_o^a\) is a dimension-only factor \(1+O(1/n)\)
by Stirling's formula. Together with the exact conditional overlap
density above, this proves, uniformly on the regular window,
\[
 f_{oo}^{\rm unw}=\widehat f_{oo}^{\,a,a',U}(1+\varepsilon_{oo}),
 \qquad |\varepsilon_{oo}|\le C/n.
\]
No raw shell density, rank factor, or normalizing constant has
been removed from either side.

For the numerator inversion, first smooth the uncut statistic laws
only in their nonlattice coordinates on a slightly enlarged window,
and then apply the smooth projection cuts and the multiplier defining
\(g\). Let \(\widehat f_i^{\rm mod}\), for
\(i=pp,pg,gp,gg\), denote the corresponding formal pair function
with these factors included. Uniform inversion gives at each pair's
own mixed point
\[
 f_i^{\rm sm}(w)
 =J_\Phi\widehat f_i^{\rm mod}(\Phi(w))(1+\varepsilon_i(w)),
 \qquad |\varepsilon_i(w)|\le Cq/N.
\]
The factors are applied after smoothing, so this identity makes no
relative-error assertion at a zero of a convolved cutoff. Commuting
those factors through a mollifier of radius \(\delta\) instead
changes the tested integral by at most \(\delta n^C\) times the
positive mass on the enlarged window. The multiplier's logarithmic
derivatives and the cutoff derivatives have polynomial bounds there. The
positive-mass estimate below and a sufficiently high inverse-polynomial
choice of \(\delta\) make this separate error negligible.

Write \(\widehat R_i=\widehat f_i^{\rm mod}/\widehat f_{oo}\)
in this chart, and let \(L_i\) be the finite smooth individual
expansion used above for quadrature. Thus
\(\widehat R_i=L_i+\widehat e_i\), where
\(|\widehat e_i|\le C(q/N)E_*^C e^{\zeta E_*^2}\).
Consequently the exact ratios used by the angular disintegration are
\[
 \begin{split}
 r_i(w)&=\frac{f_i^{\rm sm}(w)}
                 {J_\Phi f_{oo}^{\rm unw}(\Phi(w))}
          =L_i(w)+e_i(w),\\
 e_i(w)&=\widehat e_i(w)+\widehat R_i(w)
      \frac{\varepsilon_i(w)-\varepsilon_{oo}(\Phi(w))}
           {1+\varepsilon_{oo}(\Phi(w))},\qquad
 |e_i(w)|\le C\frac qN E_*^C e^{\zeta E_*^2}.
 \end{split}
\]
Here \(1/n\le q/N\), and the individual formal ratios have
the Gaussian-growth envelope already proved. Only the smooth
\(L_i\) undergo quadrature. Their signed sum has the leading
odd term in \eqref{prof:eq16}, with its stated envelope remainder.
These are absolute remainder bounds on each
own mixed measure, including cutoff boundaries; no regularity of
\(e_i\) is needed. Bounds \eqref{prof:eq18} integrate or sum them
against the absolute value of the signed angular average, not
against the angular average of the absolute insertion.
Thus the envelope in Lemma~\ref{prof:mixed-coordinate-transfer}
is \(\mathscr E=E_*^C e^{\zeta E_*^2}\).

Holding the length constraints exact ensures that at those points both angular averages use the same residual norms and Gram. Thus a scalar saddle error against an angular kernel is integrated using that kernel's reference signed density at the points, not by integrating absolute Wick polynomials. Mollification changes the angular test itself and therefore requires
a separate estimate. For that change we use absolute bounds: comparison of positive angular integrands at neighboring regular Grams, with polynomial derivative costs relative to base exponential weights.

Such costs times an arbitrarily small inverse polynomial are negligible here. Indeed the positive partition integrals with smooth cut to this region are polynomially bounded in reference units already by the scalar upper saddle bounds against the mass version of \eqref{prof:eq18}, summing also on the meshes; neighboring unsmoothed mass needed on a smaller region obeys the same type of bound by the angular energy comparison itself and unit mollifier mass. One can slightly enlarge regular windows in this justification and use smooth boundaries, without moving length or lattice statistics.

Cut boundaries can all use smooth tapers, neglecting the \(t\)-boundary by the preceding tails or \eqref{prof:eq18}. Lemmas~\ref{prof:common-angular-kernel} and
\ref{prof:mixed-coordinate-transfer} below state this disintegration and
quadrature step in a single explicit chart, including its Jacobian and
primitive lattice covolumes. They justify \eqref{prof:eq16} inside both
mass and signed-kernel tests without identifying the distinct atomic
and continuous statistic measures.

\par\smallskip\noindent\textit{Transfer of moments and final parameter choices.}\enspace
We now finish \eqref{prof:eq13}. All four pair laws have been related
to the same angular kernel, with their individual remainders controlled.
We can therefore square differences of the single-replica integrals.

Under Haar averaging, the prior density ratios multiply the common
angular integral at fixed statistics. When the reference prior
density and the two energy factors in that integral are combined
and divided by \(Z_*^2\), they give the tilted reference shell
kernel. This is the mixed-measure identity made explicit below.
The mass of modified \(g\), divided by \(Z_*\), tends to one in
probability: its first and second moments follow from
\eqref{prof:eq14}, the individual pair comparisons, and the tails.

The Hilbert-Schmidt second moments for the \emph{difference} of the cut \(p\) and \(g\) integrals, still with common denominator, are bounded as follows:
\begin{equation}\label{prof:eq19}
\begin{array}{c|c}
 {\rm integrand} & \mathbb E_{\rm Haar}\|\text{difference}\|_{\rm HS}^2\\ \hline
 1 & O(q/N)\\
 z & O(q L/\sqrt N)\\
 zz^T & O(q L^2/N)\\
 A_0(zz^T-C_o)A_0 & O(n q^2/N).
 \end{array}
\end{equation}
Each row follows by expressing the square as a pair integral and using
\eqref{prof:eq16}, \eqref{prof:eq18}.

The insertion determines whether
an odd-part gain is needed. For the mean the inner product kernel \(L U+a\cdot a'\) needs no odd gain. For the raw second moment its square differs from \(L^2U^2\) by at most \(L^2 E_*^4 C/\sqrt N\), so the leading odd term is harmless. For the last test use exactly the signed \(H\) bounds. All constants here may use a fixed small upper bound on \(\epsilon_1,\zeta,\epsilon_0\) satisfying the preceding absorptions.

On orientations obeying \eqref{prof:eq12}, the uncut \(g_0\) physical law has centered covariance at most \((1+o(1))K^{-1}\) and raw second moment at most \(C K^{-1}\), exactly by the shell argument of \eqref{prof:eq2}. Here displacement of the standardized length target is bounded rather than vanishing, so use centered conditional Gaussian variance in that argument (joint Gaussian limit with any white projection, variance can only decrease). Projection and length density bounds needed for uniform moments follow after any bounded codimension compression by \eqref{prof:eq10}.

The modification defining \(g\) changes these estimates by
\(o(1)\) in covariance units. Indeed \(W\) has uniform sub-Gaussian
tails, \eqref{prof:eq14} gives a factor tending to one with
arbitrarily small Gaussian growth, and the removed tail begins at
order \(\sqrt N\). Cauchy--Schwarz applies uniformly to each
projection moment test.

We spell out the normalization and centering. Write
\[
 I_i(F)=Z_*^{-1}\int F(z)e^{z^T(J-bI)z/2}\,i(dz),
 \qquad M_i=I_i(1),\qquad i=p,g,
\]
where both priors include their cuts and \(g\) includes its
projection multiplier. Since \(L/\sqrt N=q^{-2}\), the four
root-mean-square errors in \eqref{prof:eq19} are, in order,
\[
 O\bigl(\sqrt{q/N}\bigr),\qquad O(q^{-1/2}),\qquad
 O(q^{-3/2}),\qquad O(q^{-1/2}).
\]
Choose the failure probabilities in Markov's inequality in advance,
with their sum and the preceding orientation losses below \(1/4\).
For small \(q\) and large \(N\), the resulting orientations have
\(M_p,M_g\) bounded away from zero and close to one. For any matrix
or vector test \(F\),
\[
 \frac{I_p(F)}{M_p}-\frac{I_g(F)}{M_g}
 =\frac{I_p(F)-I_g(F)}{M_p}
   +\frac{M_g-M_p}{M_p}\frac{I_g(F)}{M_g}.
\]
The normalized modified Gaussian \(g\) has raw second moment
\(O(q^{-2})\), so its mean is \(O(q^{-1})\). For \(F=z\) and
\(F=zz^T\), the last display and
the first three errors show that the normalized mean difference is
\(O(q^{-1/2})+o(q^{-1})\), and the normalized raw second-moment
difference is \(O(q^{-3/2})+o(q^{-2})\). Subtracting the two mean
outer products costs at most \(O(q^{-3/2})+o(q^{-2})\).
Consequently the cut spin covariance is bounded above by
\((1+o(1))K^{-1}+o(q^{-2})I\).

For the weighted estimates choose the loss in
\eqref{prof:eq10} separately, denoting it by \(\eta_{\rm mat}\),
with \(2\eta_{\rm mat}<\gamma\). The normalized modified Gaussian
\(g\) has raw weighted second moment \(O(q^{-2\eta_{\rm mat}})\).
Use the last row of \eqref{prof:eq19} in the normalization identity
with \(F=A_0(zz^T-C_o)A_0\). The subtraction is a constant under
each normalized law, so the difference of raw weighted second
moments is
\[
 O(q^{-1/2})+o(q^{-2\eta_{\rm mat}}+1)=o(q^{-\gamma}).
\]
Here \(\gamma>1/2\) is exactly what makes the first error smaller
than the target. The cut spin raw weighted second moment is
therefore \(o(q^{-\gamma})\). It bounds the weighted covariance;
Jensen's inequality bounds the squared weighted mean by the same
quantity. In the unweighted case the raw second moment is
\(O(q^{-2})\), so the mean is \(O(q^{-1})=o(q^{-1-\eta})\) for
the independently prescribed \(\eta>0\).

Finally restore the removed spin mass. If it is at most \(n^{-D}\),
then \(\|X-m\|\le2\sqrt n\) bounds the changes in normalized means
and raw second moments by \(Cn^{1/2-D}\) and \(Cn^{1-D}\),
respectively; the same bounds, up to constants, hold after \(A_0\)
since \(\|A_0\|\le5\). Choose \(D\) large. The sharp unweighted
coefficient follows by choosing the precision, shell, transfer, and
\(q/r\) tolerances so that
\[
 (1+o(1))(1-\epsilon_{\rm mat})^{-1}s/q^2\le1+\epsilon.
\]
This proves \eqref{prof:eq13} with all the prescribed small
constants.

The same event also provides the mass anchor needed in the next
section. We may require \(M_p\ge1/2\), and positivity of the omitted
integral gives
\begin{equation}\label{prof:mass-anchor}
 \int e^{z^T(J-bI)z/2}\,p_{\rm full}(dz)\ge Z_*/2.
\end{equation}
It holds simultaneously with the norm bounds on the retained
orientation set. On any fixed compact observation-time interval
bounded away from zero, contained in the established small-time regime,
and on which \(r\) is bounded below, the translated posterior raw
second moment \(\E[(X-m)(X-m)^T\mid J,h]\) is bounded by a fixed
constant in operator norm on these orientations, using the covariance
and mean estimates.

All choices in this static assertion begin with coarse scalar
tolerances and matrix budgets giving fixed constants in
\eqref{prof:eq9}--\eqref{prof:eq12}, and with the orientation
failure allocations just described. Choose the overlap cutoff and
then the projection cutoff small, including for the Gaussian
envelopes, and use a sufficiently small scalar and overlap rate to
put \(q/r\) close to one and control the projection-cut loss.

Stricter diagnostics do not require increasing the coarse constants. Typical-site moment and tail cutoff exponents can be chosen for the finite derivative orders used. The upper bound on \(r\) is then sufficiently small in terms of these choices, which have not used dynamical probability/moment inputs.

\end{proof}

\subsection{Exact statistic measures and signed cancellation}

The following statements make the measure-theoretic and Gaussian parts
of the preceding comparison explicit. The smooth saddle expansions
\eqref{prof:eq14}--\eqref{prof:eq17} have already supplied the scalar
coefficients and their errors; the first two lemmas specify exactly how
these scalar quantities act on atomic and continuous priors.

\begin{lemma}[Common angular kernel and mixed statistic measures]\label{prof:common-angular-kernel}
Let $0<q<1$, $b=1-q$, $L=\sqrt{n/q}$, and $\mathsf P^T\mathsf P=nI_3$, $m=\sqrt q\mathsf P_1$, $M=n^{-1}\sum m_i$, and
$c=(1-M^2/q)^{1/2}>0$, with
$\mathsf P_2=c^{-1}(\mathbf1-(M/q)m)$.
For $z=x-m$ on $\|z\|^2+2m^Tz=nb$, set
$a=\mathsf P^Tz/\sqrt n$ and $z_\perp=z-\mathsf P a/\sqrt n$. Define primed variables
likewise, and set $U=z_\perp^Tz'_\perp/L$. Then
\[
 \rho(a)^2=nb-2\sqrt{nq}\,a_1-|a|^2,\qquad
 \Gamma(a,a',U)=
 \begin{pmatrix}\rho(a)^2&LU\\LU&\rho(a')^2\end{pmatrix}.
\]
On $\Gamma\succ0$, the statistic $(a,a',U)$ is a complete
invariant for simultaneous proper rotations of
$\operatorname{ran}(\mathsf P)^\perp$, provided $n-3\ge3$.
Thus every integrable physical pair insertion has one Haar angular
average, independent of the choice of prior. To specify the reference
used in the application, put \(\mathbf s=(a,a',U)\), let
\(T(z,z')=\mathbf s\), and define the exact unweighted density
\[
 f_{oo}^{\rm unw}=\frac{d\,T_\#(o\otimes o)}{d^3a\,d^3a'\,dU}
\]
on the positive Gram region. It includes both rank factors, both
raw shell densities, and \(\mathfrak c_o^2\). It is distinct
from the formal product of single saddle functions and a
conditional saddle function.

For a pair insertion \(F\), write
\[
 \mathcal A_F(\mathbf s)=\int_{\rm Haar}
 e^{\mathcal E(gz)+\mathcal E(gz')}F(gz,gz')\,dg,
 \qquad \mathcal E(z)=z^T(J-bI)z/2,
\]
using any representative of \(\mathbf s\). All physical matrices
in \(F\) are fixed together in this vector-rotation representation.
Let \(k_F\) be the signed statistic density under two independent
normalized physical reference shells, with insertion \(F\).
Then
\[
 Z_*^{-2}f_{oo}^{\rm unw}(\mathbf s)\mathcal A_F(\mathbf s)
       =k_F(\mathbf s)
\]
for almost every regular \(\mathbf s\). This includes the mass
insertion and the centered matrix insertion of
Lemma~\ref{prof:gaussian-signed-kernel}.

There is one smooth chart for the four pair priors. Put
$S=\sum x_i$, $S'=\sum x_i'$, and $R=\sum x_ix_i'$, and use
$w=(a_1,a_3,a'_1,a'_3,S,S',R)$. Its map $\Phi:w\mapsto(a,a',U)$ is
\[
 a_2=\frac{S-nM}{\sqrt n c}-\frac{M}{\sqrt q c}a_1,\qquad
 a'_2=\frac{S'-nM}{\sqrt n c}-\frac{M}{\sqrt q c}a'_1,
\]
\[
 U=\frac{R-nq-\sqrt{nq}(a_1+a'_1)-a\cdot a'}{L},
 \qquad |\det D\Phi|=\frac{1}{nc^2L}=:J_\Phi.
\]
For $pp$, the raw triple belongs to the affine lattice
$n(1,1,1)+\Lambda$, where
\[
 \Lambda=\operatorname{span}_{\mathbb Z}
 \{(2,0,2),(0,2,2),(0,0,4)\},\qquad
 \operatorname{covol}(\Lambda)=16.
\]
For $pg$ only $S\in n+2\mathbb Z$ is retained as a lattice
coordinate, and for $gp$ only $S'\in n+2\mathbb Z$ is retained;
these covolumes are $2$. The $gg$ measure has no lattice coordinates.

Let $dm_{pp}$ be $16$ times counting measure in its lattice
coordinates times Lebesgue measure in the remaining coordinates;
define $dm_{pg},dm_{gp}$ with factor $2$, and $dm_{gg}=dw$.
Smooth only the nonlattice coordinates, preserving the exact
translated shell constraints. For cut or multiplied priors apply
their smooth factors after this convolution; the preceding
commutator estimate charges the change from the original prior.
If $f_i^{\rm sm}$ denotes the resulting density with respect to $dm_i$, set
\[
 r_i(w)=\frac{f_i^{\rm sm}(w)}{J_\Phi f_{oo}^{\rm unw}(\Phi(w))},\qquad
 \widetilde k_F(w)=J_\Phi f_{oo}^{\rm unw}(\Phi(w))
                         \mathcal A_F(\Phi(w)).
\]
Then the smoothed pair integral is exactly
$I_i^{\rm sm}=\int r_i\widetilde k_F\,dm_i$, and
\[
 Z_*^{-2}\widetilde k_F(w)=J_\Phi k_F(\Phi(w)).
\]
For a fixed insertion we abbreviate \(\widetilde k_F\) by
\(\widetilde k\) below.
In particular, a scalar density remainder multiplies the same
signed reference kernel at each actual mixed mesh point.
\end{lemma}

\begin{proof}
Two ordered residual pairs with the same positive Gram matrix
are related by an orthogonal map. If its determinant is negative,
compose with a reflection on the orthogonal complement of the
target pair; this fixes the pair and corrects the determinant.
Haar invariance and Fubini therefore give, for every finite pair
measure $\mu$, atomic or continuous, and every integrable \(B\),
\[
 \int_G\int B(gz,gz')\,d\mu\,dg
   =\int\mathcal A_B(s)\,d(T_\#\mu)(s),
\]
where \(\mathcal A_B\) is its angular average.
For signed insertions use absolute integrability. Moving the
rotation from the interaction matrix to the pair vectors gives
precisely this identity for orbit averages in the application;
all matrix factors in the insertion move together.
For the reference prior, \(o\) is invariant under the rotations:
its isotropic Gaussian density, translated shell, and projection
rank factor are all invariant. Applying this disintegration to
\(e^{\mathcal E(z)+\mathcal E(z')}F(z,z')\,o(dz)o(dz')\),
and then using
\(Z_*^{-1}e^{\mathcal E}o=\operatorname{Law}(Z\mid T_o=0)\),
proves \(f_{oo}^{\rm unw}\mathcal A_F=Z_*^2k_F\).
For a signed insertion absolute integrability justifies Fubini;
the insertion keeps its sign inside the angular average.

The chart formulas follow by expanding $\sum x_i$ and
$x^Tx'$. Its Jacobian is triangular, with the three new diagonal
factors $1/(\sqrt n c),1/(\sqrt n c),1/L$.
For $pp$, the four sign-pattern counts are
\[
 \frac{n+S+S'+R}{4},\quad \frac{n+S-S'-R}{4},\quad
 \frac{n-S+S'-R}{4},\quad \frac{n-S-S'+R}{4}.
\]
Their integrality gives the stated lattice and determinant $16$.
The nonnegativity boundaries lie outside the central regular
window. In particular the conditional $R$ step is $4$, so three
independent step-$2$ lattices would give the wrong normalization.

The mixed lattices have only the one sign-sum coordinate.
Finally the asserted integral formula follows by multiplying and
dividing the smoothed scalar density by the positive unweighted reference
density evaluated at the same point. This is an identity in the
chosen mixed measures, not a Radon--Nikodym derivative of an
atomic statistic law with respect to a continuous law.
\end{proof}

\begin{lemma}[Separate remainders and common-volume cancellation]\label{prof:mixed-coordinate-transfer}
Suppose $r_i=L_i+e_i$, where $L_i$ extends smoothly to the common
chart, while $e_i$ need only be defined at its own mixed points.
Let $\mathscr E$ be a nonnegative envelope on the chart.
Assume $|e_i|\le\varepsilon\mathscr E$ and
\[
 \int\mathscr E|\widetilde k|\,dm_i\le B,\qquad
 \left|\int L_i\widetilde k\,dm_i
             -\int L_i\widetilde k\,dw\right|\le\delta_i.
\]
Then, with pair signs $+,-,-,+$,
\[
 \left|\sum_i s_iI_i^{\rm sm}
       -\int\Big(\sum_i s_iL_i\Big)\widetilde k\,dw\right|
 \le4\varepsilon B+\sum_i\delta_i.
\]
\end{lemma}
\begin{proof}
Bound the four errors separately on their own meshes, perform
quadrature on the four smooth terms, and only then add their
ordinary-volume integrals. No matching support or regularity of
the residual functions is needed. The relevant absolute value is
$|\widetilde k|$, the absolute value of the signed angular
average; the insertion is not replaced by its absolute value.
\end{proof}

For clarity, centered-cell quadrature supplies the displayed
$\delta_i$ from two derivatives. If $C_\Lambda$ is a centered
fundamental parallelepiped of a fixed-dimensional lattice, Taylor's
formula gives
\[
 \left|\operatorname{covol}(\Lambda)\sum_{\ell}F(\ell)
                   -\int F\right|
 \le\frac12\sum_{\ell}\int_{C_\Lambda}|u|^2
       \sup_{0\le t\le1}\|\nabla^2F(\ell+tu)\|\,du.
\]
The linear term integrates to zero. In the application,
$c^2=1-M^2/q=1-O(q)$ is bounded below. Raw lattice steps then
change $a,a'$ by $O(n^{-1/2})$ and $U$ by $O(L^{-1})$ on the
projection cut. Hold the nonlattice coordinates
$(a_1,a_3,a'_1,a'_3)$ fixed when differentiating in the raw
lattice coordinates. The only quadratic term of $\Phi$ in those
directions has
$\partial_S\partial_{S'}U=-1/(nc^2L)$.
Thus a summable polynomial--Gaussian
envelope for two derivatives in $(a,a',U)$ costs $O(1/n)$ times
its envelope scale. The covolumes in this estimate are exactly
those already present in the local Fourier inversion.

Mollification is a different error and is estimated before using
signed scalar remainders. For a probability mollifier of radius
$\delta$ in the nonlattice chart coordinates,
\[
 \int\mathcal A\,d(\nu*\kappa_\delta)-\int\mathcal A\,d\nu
 =\int\!\int[\mathcal A(w+h)-\mathcal A(w)]
                         \kappa_\delta(h)\,dh\,d\nu(w).
\]
On the enlarged regular Gram region choose a smooth Cholesky
representative and impose
$|\nabla\log\mathcal A_1|\le n^C$,
$|\nabla\mathcal A_H|\le n^{C'}\mathcal A_1$, where
$\mathcal A_1>0$ is the positive angular energy integral.

These inequalities follow by differentiating the representative:
Gram eigenvalues are of order $n$, all physical parameters and
the polynomial insertion are polynomially bounded, and the
exponential derivative is its value times a polynomial. They
preserve the translated lengths by recomputing $\rho(a)$.
A smoothed positive-mass bound then bounds the unsmoothed mass
within $e^{n^C\delta}$, so the signed smoothing error is at most
$\delta n^{C'}e^{n^C\delta}$ times that positive mass.
Taking an arbitrarily high inverse-polynomial smoothing scale
makes this negligible. Smooth enlarged cuts handle boundaries.
No Gaussian length coordinate and no retained lattice coordinate
is convolved in this argument.

\subsection{The centered Gaussian matrix kernel}



\begin{lemma}[Centered Gaussian kernel and its odd part]\label{prof:gaussian-signed-kernel}
Fix a derivative order $h\ge0$. There is a threshold $N_0=N_0(h)$
for the following statement. Let $0<q\le q_0<1$,
$N=nq^3\ge N_0$, and $L=\sqrt{n/q}$.
Let $C$ be positive definite, and let $E$ have a fixed number $d$ of
orthonormal columns. Put $\Pi=I-EE^T$. Assume
\[
 c_0I\preceq C,\quad \|C\|\le C_0q^{-2},\quad
 \operatorname{Tr}C\le C_0n,\quad \operatorname{Tr}C^2\le C_0n/q,
\]
and that at least $c_0N$ eigenvalues belong to
$[c_0q^{-2},C_0q^{-2}]$. Suppose $m\in\operatorname{ran}E$ and
$\|C^{1/2}m\|/L\le C_0$. Let $A$ be symmetric, and assume
$D=C^{1/2}A^2C^{1/2}$ has norm at most $C_0$.
For independent standard normal vectors $\eta_1,\eta_2$, set
\[
 z_i=C^{1/2}\eta_i,\quad
 T_i=(\|z_i\|^2+2m^Tz_i-nb)/L,\quad
 a_i=E^Tz_i,\quad U=z_1^T\Pi z_2/L.
\]
Assume the density of each $T_i$ at zero belongs to $[c_0,C_0]$.
Write $k_1(a_1,a_2,U)$ for the density conditional on $T_1=T_2=0$,
and $k_H$ for the signed density with insertion
\[
 H=(\eta_1^TD\eta_2)^2-\eta_1^TD^2\eta_1
       -\eta_2^TD^2\eta_2+\operatorname{Tr}D^2.
\]
Put $\alpha_i=(E^TCE)^{-1/2}a_i$ and $J_E=\det(E^TCE)^{-1}$.
For every multi-index $\beta$ with $|\beta|\le h$, on $|U|\le\epsilon\sqrt N$,
with sufficiently small fixed $\epsilon$, one has
\[
 |\partial^\beta k_1|\le C_\beta P_\beta
 e^{-c(|\alpha_1|^2+|\alpha_2|^2+U^2)}J_E,\qquad
 |\partial^\beta k_H|\le C_\beta nq P_\beta
 e^{-c(|\alpha_1|^2+|\alpha_2|^2+U^2)}J_E,
\]
where $P_\beta$ is a fixed polynomial in
$1+|\alpha_1|+|\alpha_2|+|U|$.
After integrating both projection variables, the resulting densities
$\overline k_1,\overline k_H$ obey
\[
 |\overline k_1(U)-\overline k_1(-U)|
 \le CN^{-1/2}P(U)e^{-cU^2},\qquad
 |\overline k_H(U)-\overline k_H(-U)|
 \le CnqN^{-1/2}P(U)e^{-cU^2}.
\]
The same bounds hold for derivatives of order at most $h$.
\end{lemma}

\begin{proof}
Both $B=C/L$ and $B_\Pi=C^{1/2}\Pi C^{1/2}/L$ have operator norm
$O(N^{-1/2})$, Hilbert--Schmidt norm $O(1)$, and nuclear norm
$O(\sqrt{nq})$. Their difference has rank at most $d$.
Conditioning on $a_i$ amounts to writing
\[
 \eta_i=S\alpha_i+\xi_i,\qquad
 S=C^{1/2}E(E^TCE)^{-1/2},\qquad S^TS=I,
\]
where $\xi_i$ is standard Gaussian on $S^\perp$.

Take Fourier transforms in the two lengths and in $U$, allowing a
real cross tilt $u$ with $|u|\le c\sqrt N$. The precision then has
real part bounded above and below by fixed positive constants.
Let $\Sigma,\mu$ denote the complex Gaussian covariance and mean.
The resolvent identity, the preceding three norm bounds, and the
fixed-rank projection correction give
\[
 \|\Sigma-I\|_{\rm HS}\le CP,\qquad
 \|\Sigma-I\|_1\le C\sqrt{nq}\,P,\qquad \|\mu\|\le CP,
\]
where $P$ is a polynomial in the frequencies, $|u|$, and
$1+|\alpha_1|+|\alpha_2|$. The length linear coefficient is
$C^{1/2}m/L$; projection conditioning adds coefficients of norm
at most $C(|\alpha_1|+|\alpha_2|)/\sqrt N$.

Write $\Sigma_{11}=I+\Delta_1$, $\Sigma_{22}=I+\Delta_2$,
$R=\Sigma_{12}$, and $s=\mu_1^TD\mu_2$. Wick's formula is exactly
\begin{align*}
 \mathbb E_{\Sigma,\mu}H={}&
 \operatorname{Tr}(D\Delta_2D\Delta_1)
 +\operatorname{Tr}(DR^TDR^T)+(\operatorname{Tr}(DR^T))^2\\
 &+2s\operatorname{Tr}(DR^T)+s^2
 +\mu_1^TD\Delta_2D\mu_1+\mu_2^TD\Delta_1D\mu_2
 +2\mu_1^TDR^TD\mu_2.
\end{align*}
All zero-mean terms except the trace square cost a polynomial by
Hilbert--Schmidt bounds. The exceptional term satisfies
\[
 |\operatorname{Tr}(DR^T)|^2\le\|D\|^2\|R\|_1^2\le CnqP.
\]
The mean terms cost $C\sqrt{nq}P$. In particular the terms
$\operatorname{Tr}D^2$ and $\operatorname{Tr}(D^2\Delta_i)$
have canceled identically. This is an estimate for the signed
Gaussian expectation, not for the absolute insertion.

For real Fourier variable $v\in\mathbb R^3$, determinant damping is
\[
 |\det(I-2iQ)|^{-1/2}\le C(1+c|v|^2/N)^{-c'N}.
\]
Before fixed-rank changes, $Q$ is the tensor product of $C/L$ with
the symmetric two-by-two frequency matrix. The latter has an
eigenvalue of magnitude at least a constant times $|v|$. The spread
assumption supplies $cN$ singular values of size $c|v|/\sqrt N$;
interlacing loses only a fixed number. A real cross tilt conjugates
the Fourier quadratic matrix by bounded invertible operators, so
the same bound holds with changed constants.

Every fixed frequency
polynomial is integrable against this bound, uniformly for large
$N$. Fourier inversion therefore proves the density bounds and
their fixed derivatives. To obtain the Gaussian tail in $U$, shift
the cross contour by $u=c_1U$. Its centered determinant cost is
$O(u^2)$, since the cross tilt has trace zero and bounded
Hilbert--Schmidt norm. The projected mean cost is at most
$C\epsilon(|\alpha_1|^2+|\alpha_2|^2)$ on the stated window.
Choose $c_1$ and then $\epsilon$ small. The factor $e^{-c_1U^2}$
and the projection Gaussian density absorb these costs. Division
by the bounded shell denominators completes these bounds.

The parity estimate uses the unconditioned projection variables.
Integrate them out before Fourier inversion, so that reversing the
cross frequency can be compared at the level of the Gaussian formula.
Reversing $U$ reverses both the cross Fourier frequency and the
real cross tilt. Conjugation by $\operatorname{diag}(I,-I)$
preserves the determinant and all zero-mean contractions of $H$.
The inverse cross block satisfies the sharper operator estimate
\[
 \|\Sigma_{12}\|\le CN^{-1/2}(|v|+|u|),
\]
by the block inverse formula.

The length linear coefficients have
bounded norms times $|v|$. Their Gaussian exponential factors
therefore differ by at most $N^{-1/2}$ times a frequency polynomial.
This estimate keeps the determinant damping: the two exponents
and the segment between them have nonpositive real part, because
the linear coefficients are imaginary and the real precision is
positive. For the mass this proves the extra $N^{-1/2}$.
For $H$, its zero-mean Wick expression costs $Cnq$ times this
factor. The mean terms can be bounded directly by
$C\sqrt{nq}P$, which suffices since
$\sqrt{nq}\le nq/\sqrt N$ for $q\le1$.
Fourier inversion with the same real tilt proves the odd-part
bounds and their derivatives.
\end{proof}

\begin{lemma}[The surviving scalar coefficients]\label{prof:surviving-coefficients}
Put $v_i=1-m_i^2$ and
$\kappa_2=\langle v^2\rangle$, $\kappa_3=4\langle m^2v^2\rangle$.
For $i,j\in\{p,g_0\}$, let $F_{ij}(\beta)$ be the per-site cross
log moment after optimizing all single-replica constraints. At zero projection,
\[
 F'_{ij}(0)=0,\quad F''_{ij}(0)=\kappa_2,\quad
 F'''_{ij}(0)=\kappa_3\mathbf1_{\{i=j=p\}}.
\]
The common quadratic coefficient relative to the reference is
$c_2=(\kappa_2-b^2)/(2q^2\kappa_2b^2)$. Up to the common prefactor
$b/\sqrt{\kappa_2}$ and an error $O(qE_*^Ce^{\zeta E_*^2}/N)$,
the signed zero-projection sum is
\[
 \frac{q}{\sqrt N}e^{c_2qU^2}
 \left(-\frac{\kappa_3}{2\kappa_2^2}U+\frac{\kappa_3}{6q\kappa_2^3}U^3\right).
\]
\end{lemma}
\begin{proof}
At zero constraints both canonical site laws have mean zero and
variance $v_i$.

Their third centered moments are $-2m_iv_i$ for
signs and zero for Gaussians. The cross observable is orthogonal
to all single-replica constraints. The first derivative of the
optimizing coefficients thus vanishes, so the first three
optimized cross derivatives are its ordinary first three
cumulants. This proves the displayed identities.

Let $D_2(\beta)$ be the constraint covariance. At zero cross tilt
it is block diagonal. Each diagonal block of $D_2'(0)$ vanishes:
the corresponding third cumulant contains a centered factor from
the other replica. Hence
$\left.\partial_\beta\log\det D_2(\beta)\right|_0=0$.
The conditional saddle prefactor is proportional to
$[F''(\beta)\det D_2(\beta)/\det D_2(0)]^{-1/2}$, whose
logarithmic $t$ derivative is therefore $-F'''(0)/(2\kappa_2^2)$.
Its cubic Legendre term is $nF'''(0)t^3/(6\kappa_2^3)$.
Now use $t=qU/\sqrt N$, $nt^2=U^2/q$, and
$nt^3=U^3/\sqrt N$. Reference finite-codimension corrections
$O(t^2+1/n)$ lie within $O(q/N)$ on bounded normalized statistics.

Away from zero the Taylor estimates following \eqref{prof:eq17} organize the logarithms as
$\log R_{ij}=L_0+C_*+A_i+A'_j+B_{ij}+r_{ij}$.
Here $L_0$ is common constant/quadratic, $C_*$ is common,
$A_i,A'_j$ are additive first projection terms, $B_{ij}$ is the
odd term just computed, and the retained small terms have size
$CqE_*^C/\sqrt N$. The remainder has size
$CqE_*^Ce^{\zeta E_*^2}/N$. With $s_p=1,s_{g_0}=-1$,
$\sum s_is_j=\sum s_is_j(C_*+A_i+A'_j)=0$.
Quadratic products from exponentiation cost $q^2/N\le q/N$.

The common projection multiplier differs from one by
$N^{-1/2}E_*^Ce^{\zeta E_*^2}$, so its effect on the surviving
odd term also fits the stated remainder. The projection estimate \eqref{prof:eq14}, the derivative bounds
following \eqref{prof:eq17}, and the mixed local-limit calculation
therefore give precisely this cancellation for the smooth saddle
ratios. Lemma~\ref{prof:mixed-coordinate-transfer} passes from these
ratios to the actual mixed statistic measures.
\end{proof}

\begin{corollary}[Weighted moment comparison]\label{prof:weighted-kernel-transfer}
Use the priors, cuts, and shell-constrained reference mass $Z_*$
of the preceding mixed comparison, whose physical reference covariance
is $C_o=K_o^{-1}$. Let $A$ be symmetric and assume the hypotheses of
Lemma~\ref{prof:gaussian-signed-kernel} with $C=C_o$.
Suppose the mixed-coordinate disintegration and local saddle
estimates give the preceding expansion, with scalar remainder
$O(qE_*^Ce^{\zeta E_*^2}/N)$ at each actual mixed mesh point.
Suppose centered-cell quadrature for the smooth leading terms
costs $O(1/n)$ times the signed kernel envelope. Then, with the
common normalization by $Z_*$,
\[
 \mathbb E_{\rm Haar}\left\|
 Z_*^{-1}\int A(zz^T-C_o)A\,d(p-g)_{\rm tilted}
 \right\|_{\rm HS}^2\le Cnq^2/N.
\]
\end{corollary}
\begin{proof}
There are three errors to account for: scalar remainders, the surviving
odd term, and quadrature.

The scalar remainders cost $(q/N)(nq)$ using
the signed density bound on each mesh separately. After quadrature, the odd leading
term costs $(q/\sqrt N)(nq/\sqrt N)$. Quadrature itself costs
$(nq)/n=q\le nq^2/N$. Removing smooth projection cuts is
exponentially negligible since $|W|\le C|\alpha|$ and the cuts
start at $|W|$ of order $\sqrt N$. Mollification is separately
controlled by an absolute angular derivative bound and an
arbitrarily high inverse-polynomial smoothing scale. The scalar
saddle remainders never require placing $|H|$ inside the angular
integral.
\end{proof}
